\chapter*{Glossary}
\addcontentsline{toc}{chapter}{Glossary}
\thispagestyle{plain}

\vspace{-0.3cm}

\renewcommand{\arraystretch}{1.25}
\setlength{\tabcolsep}{7pt}

\begin{longtable}{
    >{\bfseries\raggedright\arraybackslash}p{0.27\textwidth}
    >{\raggedright\arraybackslash}p{0.67\textwidth}
}

\rowcolor{black!10}
\textbf{Term} & \textbf{Usage in this Dissertation} \\
\hline
\endfirsthead

\rowcolor{black!10}
\textbf{Term} & \textbf{Usage in this Dissertation} \\
\hline
\endhead

Co-design &
A process of collective creativity applied across the design process, in which participants contribute to exploring, developing and refining ideas and artefacts \cite{sanders2008}. In this dissertation, co-design refers particularly to the collaborative design activities undertaken with job seekers, trainers and the researcher. \\
\hline

Participatory Design (PD) &
A set of theories, practices and studies concerned with the participation of intended users and other stakeholders in design activities and decisions \cite{muller2007}. In this dissertation, Participatory Design provides the broader methodological orientation within which the co-design activities were conducted. \\
\hline

Participatory Design for Development (PD4D) &
The application and adaptation of Participatory Design approaches within development contexts. In this dissertation, the term is particularly relevant to conducting participatory technology design with job seekers and intermediary organisations within the South African context. \\
\hline

Information and Communication Technologies for Development (ICT4D) &
An academic field and field of practice concerned with the use and study of information and communication technologies in relation to social and economic development. In this dissertation, ICT4D provides the broader context for examining how ICTs may support job seeking and employment-related activities. \\
\hline

Tell--Make--Enact &
The iterative co-design framework used to describe the design activities in this research \cite{simonsen2013}. \textbf{Tell} involves participants communicating experiences, needs or ideas; \textbf{Make} involves expressing or developing ideas through representations or prototypes; and \textbf{Enact} involves engaging with, exploring or using emerging designs in ways that generate further understanding and design decisions. \\
\hline

Low-fidelity prototype &
A relatively simple and inexpensive representation of a proposed design used to explore concepts, requirements and interactions before substantial implementation. In this dissertation, early paper-based and other simplified representations are treated as low-fidelity prototypes. \\
\hline

High-fidelity prototype &
A more developed representation of an intended system that more closely approximates the appearance, interaction or functionality of the eventual artefact. \\
\hline

Artefact &
The technological intervention developed through the co-design process. In this dissertation, this primarily refers to the job-searching mobisite developed and subsequently piloted with participants. \\
\hline

Job seeker &
An individual seeking employment. In this dissertation, the term primarily refers to participants who were unemployed or seeking employment during the research period. \\
\hline

Job readiness &
Preparation intended to develop the knowledge, skills and capabilities needed to seek, obtain and participate in employment. \\
\hline

Job readiness classes/programmes &
Training provided to prepare participants for seeking and entering employment and, where applicable, to develop capabilities related to starting a business. \\
\hline

Low-skilled job seeker &
A term used in this dissertation for job seekers with limited formal qualifications, occupational skills or work experience. The term does not imply an absence of abilities, knowledge or other forms of expertise. \\
\hline

Intermediary &
An individual or organisation that operates between an ICT intervention and its intended users and may facilitate access, communication, support or technology use. In this dissertation, trainers and partner NGOs performed intermediary roles. \\
\hline

Ummeli &
A South African digital employment service discussed in this dissertation as an example of an ICT-based intervention intended to support job seeking \cite{gitau2012}. \\
\hline

Babajob &
An Indian digital employment service discussed in this dissertation as an example of an ICT-based intervention connecting job seekers with employers \cite{babajob2016}. \\
\hline

Employed persons &
Statistics South Africa defines employed persons as people aged 15--64 who performed at least one hour of work during the reference week, or who had a job or business but were temporarily absent from work \cite{statssa2016}. \textbf{For this dissertation, the term is used more simply to refer to participants who had work during the research period.} \\
\hline

Unemployed persons &
Statistics South Africa defines unemployed persons using criteria concerning employment status, active job search and availability for work \cite{statssa2016}. \textbf{For this dissertation, the term refers to participants who did not have work during the research period.} \\
\hline

Underemployed persons &
Statistics South Africa defines underemployed persons as employed people who were willing and available to work additional hours and whose actual working hours were below 35 hours per week during the reference period \cite{statssa2016}. \textbf{In this dissertation, the term refers to participants who had some work during the research period but had capacity and availability for additional work.} \\
\hline

Labour force &
All persons classified as either employed or unemployed \cite{statssa2016}. \\
\hline

Non-economically active persons &
Persons aged 15--64 who are classified as neither employed nor unemployed during the reference week \cite{statssa2016}. \\
\hline

Unemployment rate &
The proportion of the labour force that is unemployed \cite{statssa2016}. \\

\end{longtable}